\chapter{Canonical PhaseNav Dependencies and the Proof Firewall}
\label{ch:phasenav-dependency-canon}

The preceding chapters develop the native PhaseNav constructions internal to
this monograph.  This chapter records a separate software integration result:
the TIR--half dependency axis is now represented by a typed canonical PhaseNav
routing layer, with NOEMA carrying its provenance state.  This changes the
software status of several dependencies, but it does not change the proof
status of the Riemann Hypothesis programme.

\section{Executable parent state}

The canonical PhaseNav parent for this integration is commit
\texttt{54f65f2ca7d35cdd98f0ab8984cc1a8d74444a96}, hard-canon version
\texttt{1.2.0}, dependency-layer version \texttt{0.7.0}.  The corresponding
NOEMA dependency/provenance parent is commit
\texttt{42a0a8916e81ca27f2213bf0f28538f046c2e89a}.

These hashes identify executable software states.  They are provenance anchors,
not mathematical axioms and not empirical evidence.

\section{Intrinsic and routing dimensions}

The recovered TIR coefficient state has the intrinsic form
\[
  (h,a,b,c)\in\mathbb{Z}^{4},
\]
with generator
\[
  G(h,a,b,c)
  =\frac{h}{2}+a\kappa+b\frac{\kappa}{L_3}
   +c\frac{\kappa^2}{2}.
\]
The PhaseNav representation used for routing and holonomic computation is a
separate point of $T^{36}$.  Consequently
\[
  \dim(\text{coefficient state})=4,
  \qquad
  \dim(\text{routing envelope})=36,
\]
and there is no identification of the two spaces.  In particular, the 36D
vector is forbidden as a hidden source of coefficient values.

\section{Refined assignment ledger}

The earlier bridge correctly kept the role/orientation problem open, but treated
its internal software structure too coarsely.  The v0.7 ledger separates four
levels:

\begin{description}
  \item[Slot-role routing -- structural PASS.]
  The four slots $h,a,b,c$ have an explicit structural routing grammar.

  \item[Directed-orbit orientation -- implemented project rule.]
  The canonical runtime contains an executable directed-orbit convention.  Its
  implementation does not establish that the convention is a physical law.

  \item[Relational-gradient orientation -- implemented source operator,
  candidate TIR binding.]
  The operator exists and can be executed, but its assignment to the physical
  coefficient slots has not been proved prospectively.

  \item[Physical TIR slot binding -- OPEN.]
  The remaining target is a prospective map from coefficient-free atomic or
  orbital state to sign, orientation, and the recovered coefficient tuple.
\end{description}

Measured masses, measured Yukawa couplings, and coefficient-enriched PhaseNav
states are disallowed as parents of this derivation.  This is a no-circularity
condition rather than a numerical tolerance.

\section{NOEMA authority boundary}

NOEMA receives the PhaseNav commit identity, the typed epistemic statuses, the
state-space dimensions, and the associated receipts.  The contract is
read-only with respect to the scientific claim ledger: a memory write, recall,
or consolidation event cannot promote a claim.  NOEMA therefore acts here as a
trajectory/provenance store, not as a source of mathematical truth.

\section{Proof firewall}

The principal open bridge remains unchanged.  In the notation of the claim
ledger:

\begin{description}
  \item[SOH-C004 -- OPEN.] Every non-trivial zero of $\xi$ must be shown to
  close in the canonical self-dual/native PhaseNav shell.

  \item[SOH-C005 -- OPEN.] Dense-family arithmetic Weil positivity,
  regularization of the complete form, and the implication from null structure
  to native closure remain to be established.
\end{description}

The new dependency layer does not supply either theorem.  In particular,
vectorization, holonomy computation, a finite positive matrix sample, a TIR
phenomenological relation, or a NOEMA receipt cannot promote these open claims.

\section{Cross-project falsification state}

For provenance, the bridge also carries selected TIR statuses.  The common
likelihood with uncertainties and covariance remains open, the retrospective
Collatz stopping-orientation pattern still requires a prospective test, and the
current TIR neutron-EDM relation remains a physical failure.  Their inclusion is
intended to prevent cross-project dependency drift, not to import their
phenomenology as evidence for the half-axis programme.

\section{What has actually been gained}

The v0.7 integration closes an architectural ambiguity.  A consumer can now
identify which pieces are exact anchors, project-derived algebra, implemented
runtime structure, open physical bindings, retrospective candidates, or failed
physical claims before following an edge through PhaseNav or NOEMA.  This makes
cross-project execution stricter without weakening the claim ledger.

The mathematical debt is therefore more sharply localized rather than declared
paid: the next promotion step remains a theorem or reproducible construction
closing the open positivity/native-closure bridge, not another layer of software
routing.
